$N=420$ works with divisors $4,6,10,14,15,21,35$. For smaller $N$, consider its number of distinct prime factors. With one prime, all divisors form a chain. With two primes of exponents $a\ge b$, partition divisors into $b+1$ chains by fixing the smaller prime\textquotesingle s exponent; seven incomparable divisors would force $b\ge6$, hence $N\ge6^6$. With three primes, sort exponents $a\ge b\ge c$. If $c\ge2$, then $N\ge2^2 3^2 5^2=900$; if $b\ge3$, then $N\ge2^3 3^3 5=1080$. Otherwise $c=1,b\le2$, and fixing those two exponents partitions divisors into at most six chains. With four primes, $N<420$ forces all exponents to be one; divisors correspond to subsets of four elements, whose largest antichain has size six (the elementary four-element case of Sperner\textquotesingle s inequality). Five primes already have product at least 2310. Thus $\boxed{420}$ is minimal.
